\documentclass[11pt]{article}
\usepackage{amsmath,amssymb}
\usepackage[margin=1in]{geometry}
\usepackage{listings}
\usepackage{xcolor}
\usepackage{array}
\usepackage{booktabs}
\usepackage{hyperref}

\definecolor{codebg}{gray}{0.96}
\lstset{
  language=Python,
  basicstyle=\ttfamily\footnotesize,
  backgroundcolor=\color{codebg},
  keywordstyle=\color{blue!70!black},
  commentstyle=\color{green!45!black},
  stringstyle=\color{red!60!black},
  numbers=left, numberstyle=\tiny\color{gray}, numbersep=6pt,
  breaklines=true, frame=single, framesep=4pt, columns=fullflexible,
  showstringspaces=false,
}

\title{Two $6k\pm1$ Factoring Methods: the Wheel Scan and the Small-Factor Sieve\\
\large A comparison}
\author{Shaimaa Soltan}
\date{August 2026}

\begin{document}
\maketitle

\begin{abstract}
Both methods rest on the single fact that, after peeling $2$ and $3$, every prime
factor is of the form $6k\pm1$. The first, \texttt{wheel\_scan}, is $6k\pm1$
\emph{trial division} with early exit: it returns the smallest factor of a single
$N$. The second builds a reusable list $P$ of small factors once, then screens many
$N$ against it. We give both, then compare them on completeness, cost, memory, and
behaviour on large semiprimes.
\end{abstract}

\section{Method 1 --- the wheel scan}

Peel $2$ and $3$; then scan the candidates $M=6m-1$ and $M=6m+1$ for
$m=1,2,3,\dots$ up to $\sqrt N$, returning at the first $M$ with $N\bmod M=0$.
Covering \emph{both} residue classes $\{6m-1,\,6m+1\}$ is essential --- a factor may
lie in either (e.g.\ $55=5\cdot11$ has $N\equiv1\bmod6$ but both primes
$\equiv5$).

\begin{lstlisting}
def wheel_scan(N):
    """Scan M = 6m-1 and 6m+1 for m = 1,2,3,...; EXIT at the first M with N%M==0.
    Peels 2 and 3 first; returns N if no factor <= sqrt(N) (N is prime)."""
    for p in (2, 3):                      # peel 2 and 3
        if N % p == 0:
            return p
    lim = isqrt(N)
    m = 1
    while 6 * m - 1 <= lim:
        for M in (6 * m - 1, 6 * m + 1):  # both classes: 5,7 / 11,13 / 17,19 / ...
            if M <= lim and N % M == 0:
                return M                  # early exit at the first m giving mod 0
        m += 1
    return N                              # prime
\end{lstlisting}

\noindent\textbf{Properties.} Deterministic and \emph{complete}: it returns the
smallest prime factor of $N$ if one is $\le\sqrt N$, else declares $N$ prime. Cost
$O(\sqrt N)$ in the worst case (smallest factor near $\sqrt N$); $O(1)$ memory; no
precomputation. Instant when a small factor exists (early exit).

\section{Method 2 --- the small-factor sieve}

Build a set $S$ of $6k\pm1$ numbers; sweep a range and record, in a list $P$, every
small divisor $k\in S$ of the swept numbers together with the cofactor $x/\!/k$
(\emph{integer} division). Then \texttt{N\_P} screens any $N$ against $P$.

\begin{lstlisting}
L = {}
for x in range(1, 10000, 2):
    if x % 6 == 5:                 # what B==8.5 and D==0.5 actually select
        L[x] = x + 2               # the pair (6k-1, 6k+1)

L2_Y = [k % v for v in L.values() for k in L.keys()
        if k != v and v % k == 0 and k % v not in L.values()]
L2_X = [v % k for k in L.keys() for v in L.values()
        if k != v and k % v == 0 and v % k not in L.keys()]
S = set(L2_X) | set(L2_Y)

P = []
for x in range(10001, 100000, 6):
    if x % 2 == 0 or x % 3 == 0:
        continue
    for k in S:
        if (x % k == 0) and k != 1 and k != x:
            if k not in P:      P.append(k)
            if x // k not in P: P.append(x // k)   # //  (integer!) -- not /

def N_P(N):
    L = [k for k in P if N % k == 0]
    print(f"Factors of {N}: {L}")
\end{lstlisting}

\noindent\textbf{Properties.} After a one-time $O(\text{range}\times|S|)$ build,
each query is only $O(|P|)$ --- fast for screening \emph{many} numbers. But it is
\emph{partial}: it finds only factors that live in $P$ (small ones), and it lists
\emph{composite} divisors too (e.g.\ $161=7\cdot23$). It also skips multiples of $3$
while building $P$, so it cannot report the factor $3$ itself.

\paragraph{A subtle bug worth flagging.} The original used
\texttt{P.append(x/k)} (true division), putting \emph{floats} like \texttt{13763.0}
in $P$. For a large integer $N$, \texttt{N \% 13763.0} first rounds $N$ to $\sim16$
digits, so it returns \texttt{0.0} by accident --- a \emph{false} factor
($N\bmod13763=6862\neq0$). The fix is one character: \texttt{x // k}.

\section{Comparison}

\begin{center}
\renewcommand{\arraystretch}{1.25}
\begin{tabular}{p{4.4cm} p{5.0cm} p{5.0cm}}
\toprule
& \textbf{wheel\_scan} & \textbf{sieve ($P$, N\_P)}\\
\midrule
Strategy & $6k\pm1$ trial division, early exit & precompute small-factor list, screen $N$\\
Peels $2,3$ & yes (explicit) & skips them (cannot report $3$)\\
Returns & smallest prime factor $\le\sqrt N$ & small divisors in $P$ (incl.\ composites)\\
Completeness & \textbf{complete} ($\le\sqrt N$) & \textbf{partial} (only factors in $P$)\\
Per-$N$ cost & $O(\sqrt N)$ worst case & $O(|P|)$ after build\\
Precompute & none & $O(\text{range}\times|S|)$ once\\
Memory & $O(1)$ & $O(|P|)$\\
Big semiprime, no small factor & correct but slow ($\sqrt N$) & returns \texttt{[]} (misses)\\
Best use & one-shot smallest-factor engine & screening many $N$ for small factors\\
\bottomrule
\end{tabular}
\end{center}

\paragraph{On the shared test numbers.}
\begin{center}
\renewcommand{\arraystretch}{1.25}
\begin{tabular}{p{6.6cm} p{3.2cm} p{4.4cm}}
\toprule
$N$ & \textbf{wheel\_scan} & \textbf{sieve N\_P}\\
\midrule
$4126375141919523791919279$ & $3$ (instant) & $[148021]$ (misses $3$)\\
$4126375141919523791919271$ & $7$ (instant) & $[7,23,47,161,329,1081,7567]$\\
$7334938410324762341973827534$\newline$419726354193798261$ & $2521$ (instant) & $[2521]$\\
$412637514191952379171729$ & too slow ($\sim\!4{\times}10^{8}$ steps) & $[\,]$ (no small factor)\\
\bottomrule
\end{tabular}
\end{center}

\noindent Note the two disagreements. For $4126375141919523791919279$,
\texttt{wheel\_scan} returns $3$ (it peels $3$ first), while the sieve misses $3$
entirely and reports $148021$ --- because $P$ was built skipping multiples of $3$.
For $4126375141919523791919271$ the sieve lists composite divisors
($161=7\cdot23$, $329=7\cdot47$, $1081=23\cdot47$, $7567=7\cdot23\cdot47$) alongside
the primes; \texttt{wheel\_scan} returns just the smallest prime, $7$.

\section{Conclusion}

The two methods are complementary faces of the same $6k\pm1$ idea:
\begin{itemize}
\item \texttt{wheel\_scan} is a \emph{complete, one-shot} smallest-factor finder:
correct for every $N$, instant when a small factor exists, and $O(\sqrt N)$ in the
worst case.
\item the \emph{sieve} is a \emph{reusable, partial} screen: cheap per query after
one build, ideal for testing many numbers for \emph{small} factors, but blind to
factors larger than its range (it returns \texttt{[]} on
$412637514191952379171729$, whose factors are $\sim\!2.5{\times}10^9$ and
$\sim\!1.6{\times}10^{14}$).
\end{itemize}
Neither escapes the $\sqrt N$ wall: the sieve is fast only because it does
\emph{less} (small factors only), and \texttt{wheel\_scan} is exactly
$2,3$-peeled trial division. For semiprimes with two large prime factors, both give
way to Pollard's rho / $p-1$ or ECM. The natural pipeline is: \textbf{sieve or wheel
first} (peel small factors cheaply), \textbf{then Pollard} for the large cofactor.

\end{document}
